% ============================================================================
% Fichier  : partie1/P01_C04_crypto_asymetrique.tex
% Titre    : Cryptographie (2) : Asymétrique, Signatures, ICP
% Auteur   : MINKA MI NGUIDJOÏ Thierry Emmanuel
% Version  : 1.0 -- novembre 2026
% ============================================================================
% Sources : Diffie et Hellman (1976) ; Rivest, Shamir, Adleman (1978) ;
% NIST SP 800-57 partie 1 (équivalences de sécurité) ; RFC 5280 (X.509) ;
% RFC 8446 (TLS 1.3) ; Shor (1994) ; FIPS 203, 204, 205 (août 2024) ;
% loi n° 2010/012, art. 8 et 18 (texte officiel). Exemples numériques
% classiques (RSA n = 3233 ; Diffie-Hellman p = 23), vérifiés.
% ============================================================================

\chapter{Cryptographie (2) : Asymétrique, Signatures, Infrastructures à Clés Publiques}
\label{chap:crypto-asym}
\chaptermark{Cryptographie asymétrique}

\epigraph{%
    « Nous sommes aujourd'hui au bord d'une révolution en cryptographie. »%
}{Whitfield Diffie et Martin Hellman, 1976}

\niveauChapitre{Fondamental}{Chapitre~\ref{chap:crypto-sym}~; arithmétique
modulaire élémentaire.}

\begin{boxobjectifs}
\begin{itemize}
    \item Expliquer le problème de la distribution des clés et sa solution
          asymétrique.
    \item Dérouler un échange de Diffie-Hellman et un chiffrement RSA sur de
          petits nombres.
    \item Expliquer la signature numérique et ses garanties.
    \item Décrire un certificat, une infrastructure à clés publiques et la
          vérification d'une chaîne.
    \item Décrire la négociation TLS~1.3.
    \item Situer la menace quantique et les algorithmes post-quantiques.
\end{itemize}
\end{boxobjectifs}

\bigskip

% ============================================================================
\section*{Question directrice : comment faire confiance à une clé reçue d'un inconnu~?}
% ============================================================================

Le technicien de SAHEL INFOGÉRANCE a saisi son mot de passe sur une fausse
page. Son navigateur affichait pourtant un cadenas~: la connexion était
chiffrée, mais \emph{vers l'attaquant}. Le chiffrement ne sert à rien si
l'on ne sait pas à qui appartient la clé. Ce chapitre explique comment la
cryptographie asymétrique et les certificats répondent à cette question,
et ce qu'ils ne garantissent pas.

\medskip
\textbf{Plan du chapitre.}
\begin{enumerate}
    \item La distribution des clés (section~\ref{sec:c04-distribution}).
    \item Diffie-Hellman et RSA (section~\ref{sec:c04-algos}).
    \item La signature numérique (section~\ref{sec:c04-signature}).
    \item Certificats et infrastructures à clés publiques (section~\ref{sec:c04-icp}).
    \item TLS~: tout assembler (section~\ref{sec:c04-tls}).
    \item La menace quantique (section~\ref{sec:c04-quantique}).
\end{enumerate}

\separateur

% ============================================================================
\section{La Distribution des Clés}
\label{sec:c04-distribution}
% ============================================================================

Avec le seul chiffrement symétrique, \(n\) personnes qui veulent toutes
communiquer en secret deux à deux ont besoin de \(n(n-1)/2\) clés, à
échanger par un canal déjà sûr~: 4\,950 clés pour 100 personnes. En 1976,
W.~Diffie et M.~Hellman proposent une rupture~: chaque personne a une
\emph{paire} de clés, l'une publique, l'autre privée~\cite{diffie1976}.

\begin{boxdef}
\termecle{Cryptographie asymétrique} (à clé publique)~: chaque acteur
possède une \emph{clé publique}, diffusée librement, et une \emph{clé
privée}, qu'il garde seul. Ce que l'une chiffre, seule l'autre déchiffre~;
et il est impraticable de retrouver la clé privée à partir de la clé
publique.
\end{boxdef}

Les algorithmes asymétriques sont bien plus lents que les symétriques. En
pratique, on combine les deux~: l'asymétrique sert à convenir d'une clé
symétrique de session, qui chiffre ensuite les données. C'est le
\termecle{chiffrement hybride}.

% ============================================================================
\section{Diffie-Hellman et RSA}
\label{sec:c04-algos}
% ============================================================================

\subsection{L'échange de Diffie-Hellman}

Alice et Bob conviennent publiquement d'un nombre premier \(p\) et d'un
générateur \(g\). Alice choisit un secret \(a\) et publie
\(A = g^a \bmod p\)~; Bob choisit \(b\) et publie \(B = g^b \bmod p\).
Chacun calcule \(K = B^a \bmod p = A^b \bmod p\). Un observateur voit
\(p\), \(g\), \(A\) et \(B\), mais ne sait pas retrouver \(K\) sans
résoudre le \emph{logarithme discret}.

\begin{boxinfo}
\textbf{Exemple avec de petits nombres.} \(p = 23\), \(g = 5\), \(a = 6\),
\(b = 15\). Alice publie \(A = 5^6 \bmod 23 = 8\)~; Bob publie
\(B = 5^{15} \bmod 23 = 19\). Alice calcule \(19^6 \bmod 23 = 2\)~; Bob
calcule \(8^{15} \bmod 23 = 2\). Ils partagent \(K = 2\). En pratique,
\(p\) a plus de 2\,000 bits, ou l'on utilise des courbes elliptiques.
\end{boxinfo}

\begin{boxalert}
\textbf{Limite.} Diffie-Hellman seul ne dit pas \emph{avec qui} on
échange. Un attaquant placé au milieu fait un échange avec Alice et un
autre avec Bob. Il faut authentifier les valeurs échangées~: c'est le
rôle de la signature et des certificats.
\end{boxalert}

\subsection{RSA}

En 1978, R.~Rivest, A.~Shamir et L.~Adleman publient le premier système
qui permet à la fois de chiffrer et de signer~\cite{rivest1978}. Sa
sécurité repose sur la difficulté de factoriser le produit de deux grands
nombres premiers.

\begin{tableaucours}{L{4.4cm}L{7.8cm}}{%
    \tableauHeader{Étape} &
    \tableauHeader{Exemple (petits nombres)}
}
Choisir deux premiers & \(p = 61\), \(q = 53\) \\
Module & \(n = pq = 3\,233\) \\
Indicatrice & \(\varphi(n) = (p-1)(q-1) = 3\,120\) \\
Exposant public & \(e = 17\) (premier avec 3\,120) \\
Exposant privé & \(d = 2\,753\), car \(17 \times 2\,753 = 46\,801 = 15
\times 3\,120 + 1\) \\
Chiffrer \(m = 65\) & \(c = 65^{17} \bmod 3\,233 = 2\,790\) \\
Déchiffrer & \(2\,790^{2\,753} \bmod 3\,233 = 65\) \\
\end{tableaucours}
\vspace{-0.8em}
\begin{center}{\small\itshape Tableau~4.1. RSA pas à pas}\end{center}

\subsection{Tailles de clés équivalentes}

\begin{tableaucours}{L{3.2cm}L{3.0cm}L{3.0cm}L{3.0cm}}{%
    \tableauHeader{Niveau (bits)} &
    \tableauHeader{Symétrique} &
    \tableauHeader{RSA, DH} &
    \tableauHeader{Courbes elliptiques}
}
112 & 3DES & 2\,048 & 224 \\
128 & AES-128 & 3\,072 & 256 \\
192 & AES-192 & 7\,680 & 384 \\
256 & AES-256 & 15\,360 & 512 \\
\end{tableaucours}
\vspace{-0.8em}
\begin{center}{\small\itshape Tableau~4.2. Équivalences de sécurité (d'après le
NIST SP~800-57)~\cite{nist80057}}\end{center}

% ============================================================================
\section{La Signature Numérique}
\label{sec:c04-signature}
% ============================================================================

\begin{boxdef}
\termecle{Signature numérique}~: valeur calculée sur l'empreinte d'un
document avec la clé \emph{privée} du signataire, et vérifiable par tous
avec sa clé \emph{publique}. Elle garantit l'intégrité du document,
l'authenticité de son origine et, parce que seul le signataire détient la
clé privée, la non-répudiation.
\end{boxdef}

\begin{boxjuris}
\textbf{\loicam{n°~2010/012}, article~18} (texte officiel)~:
\textit{« Une signature électronique avancée doit remplir les conditions
ci-après~: les données afférentes à la création de la signature sont liées
exclusivement au signataire et sont sous son contrôle exclusif~; toute
modification à elle apportée, est facilement décelable~; elle est créée au
moyen d'un dispositif sécurisé dont les caractéristiques techniques sont
fixées par un texte du Ministre chargé des Télécommunications~; le
certificat utilisé pour la génération de la signature est un certificat
qualifié. »}

\medskip
\textbf{Analyse.} Chaque condition a sa traduction technique~: clé privée
sous contrôle exclusif (stockée sur une carte ou un module matériel),
intégrité (signature sur l'empreinte), dispositif sécurisé, certificat
qualifié délivré par une autorité autorisée. L'article~17 donne à cette
signature avancée la même valeur que la signature manuscrite.
\end{boxjuris}

% ============================================================================
\section{Certificats et Infrastructures à Clés Publiques}
\label{sec:c04-icp}
% ============================================================================

Une clé publique n'est qu'un nombre~: rien ne dit à qui elle appartient.
Un \termecle{certificat} est un document signé par une \emph{autorité de
certification}, qui lie une clé publique à une identité, pour une durée
donnée. Le format standard est X.509~\cite{rfc5280}.

\begin{tableaucours}{L{3.6cm}L{8.6cm}}{%
    \tableauHeader{Champ du certificat} &
    \tableauHeader{Rôle}
}
Sujet & L'identité certifiée (un nom de domaine, une personne) \\
Clé publique & La clé liée à cette identité \\
Émetteur & L'autorité qui a signé \\
Validité & Dates de début et de fin \\
Usages & Ce que la clé peut faire (signer, chiffrer, authentifier un serveur) \\
Signature & La signature de l'émetteur sur l'ensemble \\
\end{tableaucours}
\vspace{-0.8em}
\begin{center}{\small\itshape Tableau~4.3. Les champs d'un certificat X.509}\end{center}

\begin{boxdef}
\termecle{Infrastructure à clés publiques} (ICP, ou PKI)~: ensemble des
autorités, des règles et des moyens qui émettent, publient, révoquent et
vérifient des certificats. La vérification remonte une \emph{chaîne}~:
le certificat du serveur est signé par une autorité intermédiaire, signée
à son tour par une autorité racine, à laquelle on fait confiance a priori.
\end{boxdef}

\begin{boxjuris}
\textbf{\loicam{n°~2010/012}, article~8} (texte officiel)~:
\textit{« (1) L'Agence est l'Autorité de Certification Racine. (2) L'Agence
est l'autorité de certification de l'Administration Publique. »}

\medskip
\textbf{Analyse.} Au Cameroun, l'ANTIC est la racine de la confiance
pour les certificats qualifiés. L'activité de certification est soumise à
autorisation préalable (art.~10).
\end{boxjuris}

\begin{boxprocedure}
\textbf{Vérifier un certificat}
\begin{enumerate}
    \item Le nom du sujet correspond-il au site ou à la personne attendus~?
    \item Le certificat est-il dans sa période de validité~?
    \item Chaque signature de la chaîne est-elle valide, jusqu'à une racine
          de confiance~?
    \item Le certificat a-t-il été révoqué (listes de révocation, OCSP)~?
    \item Les usages autorisés couvrent-ils l'usage qui en est fait~?
\end{enumerate}
\end{boxprocedure}

\begin{boxscenario}{Opération GAROUA. Le cadenas de la fausse page}
\textbf{Constatation.} La fausse page de connexion avait un certificat
valide, délivré automatiquement pour le domaine de l'attaquant, qui
imitait celui du fournisseur à une lettre près.

\medskip
\textbf{Analyse} (confiance élevée). Le cadenas garantit que l'on parle
bien au détenteur du domaine affiché~; il ne dit pas que ce domaine est
le bon. L'étape~1 de la procédure, vérifier le nom, aurait suffi. Un
second facteur d'authentification lié au vrai domaine (clés FIDO2) aurait
rendu la fausse page inopérante (Chapitre~\ref{chap:iam}).
\end{boxscenario}

% ============================================================================
\section{TLS : Tout Assembler}
\label{sec:c04-tls}
% ============================================================================

Le protocole TLS, dans sa version~1.3, combine tout ce qui précède~\cite{rfc8446}~:

\begin{enumerate}
    \item le client et le serveur font un échange de Diffie-Hellman sur
          courbe elliptique, qui donne une clé de session~;
    \item le serveur présente son certificat et \emph{signe} l'échange~:
          le client sait qu'il parle au détenteur du certificat~;
    \item les données sont ensuite chiffrées avec un chiffrement
          authentifié (AES-GCM ou ChaCha20-Poly1305).
\end{enumerate}

Chaque session utilise une clé éphémère~: la compromission ultérieure de
la clé privée du serveur ne permet pas de déchiffrer les sessions passées
(\emph{confidentialité persistante}).

% ============================================================================
\section{La Menace Quantique}
\label{sec:c04-quantique}
% ============================================================================

En 1994, P.~Shor montre qu'un ordinateur quantique suffisamment puissant
factoriserait et calculerait des logarithmes discrets en temps
polynomial~\cite{shor1994}~: RSA, Diffie-Hellman et les courbes
elliptiques tomberaient. Un tel ordinateur n'existe pas aujourd'hui, mais
des données chiffrées et capturées aujourd'hui pourraient être déchiffrées
demain. En août~2024, le NIST a publié trois normes post-quantiques~:
ML-KEM pour l'échange de clés (FIPS~203), ML-DSA (FIPS~204) et SLH-DSA
(FIPS~205) pour les signatures~\cite{fips203,fips204,fips205}. Les
algorithmes symétriques sont moins touchés~: doubler la taille des clés
suffit.

\begin{boxmetier}{Administrateur systèmes et réseaux}
\textbf{Tenir l'inventaire des certificats.} Un certificat expiré arrête
un service aussi sûrement qu'une attaque. L'administrateur tient la liste
des certificats, de leurs dates d'expiration et de leurs autorités, et
automatise leur renouvellement.
\end{boxmetier}

\begin{boxmetier}{RSSI}
\textbf{Préparer la migration post-quantique.} Le RSSI inventorie les
usages de la cryptographie asymétrique (VPN, TLS, signatures,
sauvegardes chiffrées) et la durée pendant laquelle chaque donnée doit
rester secrète~: ce sont les premières à migrer.
\end{boxmetier}

\separateur

\begin{boxresume}
\textbf{Ce qu'il faut retenir de ce chapitre~:}
\begin{itemize}
    \item \textbf{La cryptographie asymétrique} résout la distribution des
          clés~; on l'associe au symétrique (chiffrement hybride).
    \item \textbf{Diffie-Hellman} établit une clé partagée~; \textbf{RSA}
          chiffre et signe~; tous deux doivent être authentifiés.
    \item \textbf{La signature} apporte intégrité, authenticité et
          non-répudiation~; la loi n°~2010/012 fixe les conditions de la
          signature avancée (art.~18).
    \item \textbf{Un certificat} lie une clé à une identité~; on vérifie le
          nom, la validité, la chaîne et la révocation~; l'ANTIC est la
          racine camerounaise (art.~8).
    \item \textbf{TLS~1.3} assemble ces briques~; un cadenas ne garantit pas
          que le domaine est le bon.
    \item \textbf{Les algorithmes post-quantiques} sont normalisés depuis
          2024~; la migration commence par l'inventaire.
\end{itemize}
\end{boxresume}

\begin{boxinfo}
\textbf{Points de contrôle de l'élève.} Je sais~: calculer un échange de
Diffie-Hellman et un chiffrement RSA sur de petits nombres~; expliquer
pourquoi Diffie-Hellman seul est vulnérable à l'homme du milieu~; lire
un certificat~; vérifier une chaîne~; décrire TLS~1.3 en trois étapes~;
citer les normes FIPS~203 à~205.
\end{boxinfo}

\bigskip

\begin{boxexo}[title={\ding{45}\quad Exercice~4.1. Nombre de clés \niveaubase}]
Combien de clés symétriques faut-il pour que 150~salariés de LiZbiZ
puissent communiquer deux à deux en secret~? Et combien de paires de clés
asymétriques~?
\end{boxexo}

\begin{boxexo}[title={\ding{45}\quad Exercice~4.2. Diffie-Hellman \niveaubase}]
Avec \(p = 23\) et \(g = 5\), Alice choisit \(a = 4\) et Bob \(b = 3\).
Calculez \(A\), \(B\) et la clé partagée.
\end{boxexo}

\begin{boxexo}[title={\ding{45}\quad Exercice~4.3. RSA \niveaumoyen}]
Avec \(p = 5\), \(q = 11\) et \(e = 3\), calculez \(n\), \(\varphi(n)\) et
\(d\)~; chiffrez puis déchiffrez \(m = 7\).
\end{boxexo}

\begin{boxexo}[title={\ding{45}\quad Exercice~4.4. Quelle clé~? \niveaubase}]
Pour chaque opération, dites quelle clé utilise Alice~: chiffrer un
message pour Bob~; signer un message~; vérifier la signature de Bob~;
déchiffrer un message reçu.
\end{boxexo}

\begin{boxexo}[title={\ding{45}\quad Exercice~4.5. Lire un certificat \niveaumoyen}]
Un certificat a pour sujet \texttt{portail.lizbiz.cm}, pour émetteur une
autorité intermédiaire, une validité du 1\textsuperscript{er}~janvier au
31~mars~2027, et l'usage « authentification de serveur ». Le portail est
consulté le 2~avril à l'adresse \texttt{www.lizbiz.cm}. Listez les
problèmes.
\end{boxexo}

\begin{boxexo}[title={\ding{45}\quad Exercice~4.6. L'homme du milieu \niveaumoyen}]
Décrivez l'attaque de l'homme du milieu contre un Diffie-Hellman non
authentifié, puis expliquez comment TLS~1.3 l'empêche.
\end{boxexo}

\begin{boxexo}[title={\ding{45}\quad Exercice~4.7. Signature avancée \niveaumoyen}]
LiZbiZ veut signer ses bons de commande électroniquement. Pour chaque
condition de l'article~18, proposez une mesure technique.
\end{boxexo}

\begin{boxexo}[title={\ding{45}\quad Exercice~4.8. Confidentialité persistante \niveaumoyen}]
La clé privée du serveur de LiZbiZ est volée en 2027. Les sessions TLS de
2026, capturées par un attaquant, sont-elles lisibles~? Distinguez TLS~1.3
et un ancien échange de clés par RSA.
\end{boxexo}

\begin{boxexo}[title={\ding{45}\quad Exercice~4.9. Inventaire post-quantique \niveauexpert}]
Dressez, pour la BANQUE SANAGA, un inventaire des usages de la
cryptographie asymétrique, et classez-les selon l'urgence de leur
migration post-quantique. Justifiez par la durée de vie des données.
\end{boxexo}

\begin{boxexo}[title={\ding{45}\quad Exercice~4.10. La fausse page \niveauexpert}]
Proposez trois mesures, de natures différentes (humaine, technique,
cryptographique), qui auraient empêché le vol du mot de passe du
technicien, et classez-les par efficacité.
\end{boxexo}

\bigskip

\begin{boxinfo}
\textbf{Pour aller plus loin.}
\begin{itemize}
    \item W.~Diffie et M.~Hellman~\cite{diffie1976}~; R.~Rivest, A.~Shamir
          et L.~Adleman~\cite{rivest1978}.
    \item RFC~5280 (certificats) et RFC~8446 (TLS~1.3)~\cite{rfc5280,rfc8446}.
    \item FIPS~203 à~205~\cite{fips203,fips204,fips205}.
\end{itemize}
\end{boxinfo}

\begin{boxinfo}
\textbf{Transition.} Les chapitres précédents ont donné les notions et les
outils. Le Chapitre~\ref{chap:attaques} les éprouve sur l'histoire~: les
grandes attaques et ce qu'elles ont appris.
\end{boxinfo}
